\chapter{Introduzione}
\am{Nota generale. Non so se è una mia impressione ma mi sembra piccolo il font rispetto alle altre tesi che vedo di solito. Non che sia un problema, stai solo attento poi per l' eventuale stampa}\\
\am{Nota generale 2. Nella tesi ci sono pochi riferimenti (te lo dicono perchè è forse la cosa principale che solitamente le commissioni guardano). Non dico di mettere cose a caso ma magari sopratutto nella parte di background quando introduci strumenti o concetti qualche riferimento in più a paper di sicuro farà contenta la commissione }
\section{Sicurezza del Firmware nei Dispositivi Medicali}

I dispositivi embedded medicali rappresentano ad oggi una categoria particolarmente critica dal punto di
vista della sicurezza informatica. A differenza di dispositivi IT classici, questi operano in contesti in cui
un malfunzionamento o una compromissione può tradursi in rischi per la salute dei pazienti o per la loro
privacy~\cite{halperin2008pacemakers}.

La crescente connettività di questi apparati -- tramite interfacce Bluetooth, Wi-Fi, USB o altri protocolli
-- ha ampliato enormemente i vettori di attacco possibili, esponendoli a minacce che in passato erano una
prerogativa solo dei dispositivi classici quali personal computer e server. Come riportato dall'Agenzia
per la Cybersicurezza Nazionale~\cite{acn2025}, nel periodo 2023--2025 sul territorio nazionale si sono
verificati mediamente 4,7 eventi cyber malevoli al mese ai danni di strutture sanitarie, dei quali la metà
ha avuto impatti sui servizi erogati e sulla privacy dei pazienti.

Per contrastare questi attacchi sono state proposte diverse soluzioni, ma le caratteristiche intrinseche di
questi dispositivi rendono difficile l'utilizzo di strumenti di analisi pensati per sistemi standard. Ad
esempio, i sorgenti sono raramente disponibili, e i binari risultanti sono spesso privi di simboli di debug,
compilati per architetture non standard e progettati per operare in ambienti con scarse risorse hardware.

\section{Obiettivi della tesi}
L’obiettivo del progetto è quello di esplorare strumenti e metodologie per l’analisi di sicurezza di firmware embedded in dispositivi medicali, utilizzando come caso di studio il firmware dell’Holter Walk400h, un monitor cardiaco ambulatoriale rappresentativo della categoria di dispositivi medicali a bassa complessità ma con firmware non banale. 
A partire dal lavoro di reverse engineering condotto in
una precedente tesi~\cite{alberici26}, che ha permesso di ricostruire la struttura del bootloader e di
identificarne le funzioni critiche, l'obiettivo specifico di questo lavoro è valutare se strumenti di symbolic
execution quali KLEE e angr permettano di individuare e validare concretamente una vulnerabilità nel
meccanismo di aggiornamento del firmware, confrontandone al contempo l'efficacia e le limitazioni pratiche
in un contesto di firmware bare-metal.

\section{Struttura della tesi}

Il resto della tesi è organizzato come segue. Il Capitolo~2 introduce i concetti teorici alla base della
symbolic execution, le relative sfide e le principali tecniche di mitigazione e presenta gli
strumenti utilizzati nel lavoro: Ghidra, RetDec, KLEE e angr. Il Capitolo~3 descrive il dispositivo
Walk400h e il lavoro di reverse engineering del suo bootloader, individuando le funzioni critiche ai fini
dell'analisi. Il Capitolo~4 illustra la metodologia adottata, articolata nei due approcci esplorati --
RetDec~+~KLEE e Angr. Il Capitolo~5 presenta i risultati ottenuti con entrambi gli approcci e le
vulnerabilità identificate. Il Capitolo~6, infine, discute l'efficacia degli strumenti utilizzati, l'impatto di
sicurezza delle vulnerabilità individuate, propone possibili sviluppi futuri e riassume le conclusioni del lavoro.


